 \setlength{\baselineskip}{20pt}
 \begin{modify}

修改是论文写作过程中不可或缺的重要步骤,是提高论文质量的有效环节。修改的过程其实就是“去伪存真”、去糟粕取精华使论文不断“升华”的过程。

以下内容要求填写到毕业论文（设计）修改记录中：

一、毕业论文（设计）题目修改（没有可删除，序号依次递增）

原题目：

修稿后题目：

二、指导教师变更（没有可删除，序号依次递增）

原指导教师：

变更后指导教师：

三、校外毕业论文（设计）时间节点记录（没有可删除，序号依次递增）

本人于****年*月申请到******大学做毕业论文（设计），校外指导教师为：******

校内指导教师为：******。****年*月*日回到学校。
    
四、毕业论文（设计）内容重要修改记录（不可缺少）
    
包括：指导教师要求的重大修改，评阅教师要求的修改，答辩委员会提出的修改意见以及检测后的修改记录等。
    
五、毕业论文（设计）外文翻译修改记录（不可缺少）
    
六、毕业论文（设计）正式检测重复比（不可缺少）
    
修改记录正文选用模板中的样式所定义的“正文”，每段落首行缩进2字；字体：宋体，字号：小四，行距：多倍行距 1.25，间距：段前、段后均为0行。
    

~\par
\rightline{ 记录人（签字）：~~~~~~~~~~~~}\par
\rightline{ 指导教师（签字）：~~~~~~~~~~~~}


\end{modify}


